\section{Actual adelic volumes and their compact-dual normalization}
\label{sec:geometric-volumes}

\status{GV1--GV9 prove the finite and real measure conversions, with
every determinant, dyadic, discriminant, and factorial factor specified.
GV11 proves the geometric definite quotient-to-mass identity.
GV10 and GV13--GV16 use one explicitly retained classical input: the
Tamagawa-number theorem for unitary groups, with the measure convention
stated in GV10. This gives genuine signature independence and the
literal \texttt{VolumeEqMass\textsubscript{2}} and \texttt{VolTieG}
suppliers on their stated data. GV14 also uses the named
Hasse--Minkowski theorem already retained in MF4. The final stage joins
keep their exact objects. Discrete-series existence and the formal-degree
formula, as well as the remaining nonempty compatible \texttt{V4.Core}
data, are not inferred from a volume formula. No Lean build or kernel
verification is asserted.}

\paragraph{Conventions and the new classical boundary.}
Let $E/F$ be a CM quadratic extension, with $F$ totally real of degree
$f$. Unless the assertion explicitly allows ramification, assume the
literal \texttt{CMData.cmUnram} condition. DR1 and MF1 identify this
condition with an unramified quadratic algebra at every finite place,
including places above $2$. Write $\epsilon_v=1$ at a split finite
place and $\epsilon_v=-1$ at a nonsplit place; $q_v$ is the residue
cardinality of $F_v$. Absolute values on $F$ are normalized by
$|\varpi_v|_v=q_v^{-1}$ and the usual absolute value at real places.
The quadratic Hecke character of $E/F$ is denoted $\epsilon$.
For $j\ge2$ put
\[
 L_F(j,\epsilon^j)=\prod_{v<\infty}
                  (1-\epsilon_v^j q_v^{-j})^{-1}.
 \tag{GV.1}
\]
These products converge absolutely and are positive. The positive
finite value $L_F(1,\epsilon)$ and its occurrence in Tamagawa measure
are part of the precise classical theorem retained in GV10; no
conditionally convergent Euler product at $1$ is manipulated below.

A supplied \texttt{HSpace} $W$ has a nondegenerate Hermitian form,
a rational frame with Gram matrix $J$, and an actual full self-dual
lattice $L=W.L_0$. No global integral basis of $L$ is assumed. Put
$n=\dim_E W$, $G_J=U(J)$, and $\Gamma=G_J(F)$, embedded diagonally.
Thus the source's $\texttt{HA}(\mathbb A_F)$ and
$\texttt{HFinIn}(\mathbb A_F)$ are respectively
$G_J(\mathbb A_F)$ and $\Gamma$. A CM-type chooses one complex
embedding over each real place. We use the actual supplied
$\texttt{CDHaar}$, rather than an independently chosen real Haar
measure. Set
\[
 a_n=\frac{n(n+1)}2,\qquad b_n=\frac{n(n-1)}2,\qquad
 v_n=\frac{2^n\pi^{a_n}}{\prod_{j=0}^{n-1}j!},\qquad
 A_n=2^{b_n}v_n=\frac{(2\pi)^{a_n}}{\prod_{j=0}^{n-1}j!}.
 \tag{GV.2}
\]
Empty products are one, so $v_0=A_0=1$.

\paragraph{GV1: the actual adelic group and unimodularity.}
Choose an $F$-basis of $E$. Expanding each matrix entry in this basis
identifies the source tensor algebra with the usual restricted product
of $E\otimes_F F_v$, and separates the finite and real components:
\[
 G_J(\mathbb A_F)\simeq G_J(\mathbb A_{F,f})
                         \times\prod_{\tau:F\hookrightarrow\mathbb R}
                                      U(p_\tau,q_\tau).
 \tag{GV.3}
\]
The equations $g^*Jg=J$ and $g^{-1}=J^{-1}g^*J$ prove that these are
topological group identifications, not merely bijections of sets.
The finitely many exceptional coefficient denominators do not change
the restricted-product topology. Finite local matrix spaces and the
real matrix spaces are locally compact and second countable; using the
countably many finite sets of places proves the same for (GV.3).

These groups are unimodular. Indeed, after extending a local base field
to split $E/F$, the Lie algebra becomes $M_n$ and the adjoint action is
$X\mapsto gXg^{-1}$. Its determinant is
$(\det g)^n(\det g^{-1})^n=1$. This determinant identity descends to
the original field. A nonzero invariant top differential therefore
gives a Haar measure invariant under right translation as well as
left translation. At real places the same conclusion follows from
the absolute determinant on the real Lie algebra. Local products and
the restricted product are consequently also unimodular. In
particular every inner conjugation preserves Haar measure.

\paragraph{GV2: the source completion is the completed lattice.}
The symbol \texttt{AddSubgroup.closure} in \texttt{latHat} means the
algebraically generated additive subgroup. It is not a topological
closure. Nevertheless, for a full lattice $L$ it is exactly
\[
 \widehat L=\prod_{v<\infty}L_v,
 \qquad L_v=\mathcal O_{F_v}\otimes_{\mathcal O_F}L.
 \tag{GV.4}
\]
To see this, choose finitely many $\mathcal O_F$-generators $x_i$ of
$L$. At each $v$, every element of $L_v$ is $\sum_i a_{iv}x_i$ with
$a_{iv}\in\mathcal O_{F_v}$. The coefficients $a_i=(a_{iv})_v$ lie
in $\widehat{\mathcal O_F}$ and give the finite sum of pure tensors
allowed by \texttt{latHat}. Conversely every such generator belongs
to every $L_v$. This proves (GV.4) in the source's fixed rational frame.

The lattice is compact, as the image of a finite power of
$\widehat{\mathcal O_F}$. It is open: a common nonzero denominator
gives
$c\widehat{\mathcal O_E}^{n}\subseteq\widehat L
 \subseteq c^{-1}\widehat{\mathcal O_E}^{n}$.
Its stabilizer
\[
 K_L=\{g\in G_J(\mathbb A_{F,f}):g\widehat L=\widehat L\}
      =\prod_{v<\infty}K_{L_v}
 \tag{GV.5}
\]
is compact and open. In local bases a lattice stabilizer in the general
linear group is $\mathrm{GL}_n(\mathcal O_{E_v})$ (the corresponding
product at a split place); intersect with the closed unitary equations.
These descriptions agree with the standard integral model outside a
finite set. They show compactness and openness directly, and also show
that both implications in the source's \texttt{StabHat} are required.

\paragraph{GV3: existence and uniqueness of the prescribed Haar measure.}
Give $G_J(\mathbb A_{F,f})$ its unique Haar measure with
$K_L$ of mass one. At every real place choose a matrix
$A_\tau$ with
$A_\tau^*J_{p_\tau,q_\tau}A_\tau=\tau J$ and transport
$\texttt{cd.cd}(p_\tau,q_\tau)$ by $g\mapsto A_\tau gA_\tau^{-1}$.
The product through (GV.3) is a Haar measure $\mu_{W,L}$ satisfying
\[
 \mu_{W,L}\left(K_L\times\prod_\tau A_\tau^{-1}B_\tau A_\tau\right)
                    =\prod_\tau\texttt{cd.cd}(p_\tau,q_\tau)(B_\tau)
 \tag{GV.6}
\]
for all measurable $B_\tau$.

Every other signature matrix is $u_\tau A_\tau$ with
$u_\tau\in U(p_\tau,q_\tau)$. GV1 makes (GV.6) independent of these
choices. The archimedean algebra map extending a chosen embedding of
$E$ is unique: on $a\otimes x$ it is $a_\tau\tau(x)$, and the
real-coordinate idempotents isolate the prescribed $\tau$.
Thus both existential and universal parts of the source's
\texttt{AllSome} clauses hold. Conversely two Haar measures satisfying
(GV.6) are proportional by Haar uniqueness; evaluate a rectangle with
each real factor compact of positive finite measure to see that the
proportionality constant is one. This proves literal
\texttt{IsVolMeasF} existence and uniqueness, including nonconstant
signature functions, and its constant-signature \texttt{IsVolMeas}
specialization.

\paragraph{GV4: rational discreteness and genuine fundamental domains.}
The diagonal $F$ in $\mathbb A_F$ is discrete. For example, force the
finite coordinates into $\widehat{\mathcal O_F}$ and take a sufficiently
small neighborhood of zero in all real coordinates; the Minkowski
image of $\mathcal O_F$ has no nonzero point in that neighborhood.
Apply this entrywise after fixing an $F$-basis of $E$. A neighborhood
of the identity in the adelic matrix group then contains no rational
matrix except the identity. Hence $\Gamma$ is a discrete countable
subgroup of the group in GV1.

For completeness one can construct a Borel fundamental domain without
presupposing finite volume. Cover the group by countably many open
sets $U_i$ such that $\gamma U_i\cap U_i$ is empty for
$\gamma\ne1$. Such sets exist by discreteness. Put
$D_i=U_i\setminus\bigcup_{j<i,\gamma\in\Gamma}\gamma U_j$.
The union $D=\bigcup_iD_i$ meets each orbit exactly once. It is Borel,
and its translates partition the group. For any two measurable
fundamental domains $D,D'$, Tonelli and left invariance give
\[
 \mu(D)=\int_D\sum_{\gamma\in\Gamma}1_{D'}(\gamma g)\,d\mu(g)
       =\mu(D').
 \tag{GV.7}
\]
The common value is positive: otherwise countably many null translates
would cover a group carrying nonzero Haar measure. These facts also
hold for the source's almost-everywhere fundamental domains by the
same integral argument. Finiteness has not been assumed here.

\paragraph{GV5: the exact real normalization, including the powers of two.}
First use the following real coordinates on $\mathfrak u(n)$:
diagonal entries $it_i$, and, for $i<j$, entries
$X_{ij}=a_{ij}+ib_{ij}$, $X_{ji}=-a_{ij}+ib_{ij}$.
Let $dX$ be coordinate Lebesgue measure. For the map
$(X,Y)\mapsto X+iY$, a diagonal pair has real determinant of absolute
value one. An off-diagonal pair has coordinates
\[
 (a,b,c,d)\longmapsto(a-d,b+c,-a-d,b-c)
 \tag{GV.8}
\]
and determinant of absolute value four. The unique measure satisfying
the source's \texttt{IsOmegaLeb} condition is therefore
$\lambda_n=2^{b_n}dX$: its product with itself pushes to entrywise
complex Lebesgue measure exactly.

The invariant top density $dX$ has total volume $v_n$ on $U(n)$.
Here is a direct normalization proof. The first-column fibration
$U(n)\to S^{2n-1}$ has fiber $U(n-1)$. At the identity the horizontal
coordinates $t_1,a_{1j},b_{1j}$ map, with absolute determinant one,
to the ordinary tangent coordinates of the unit sphere. Splitting
the remaining coordinates gives precisely the same $dX$ on the
fiber. Invariant-measure disintegration thus yields
\[
 v_n=\operatorname{area}(S^{2n-1})v_{n-1}
     =\frac{2\pi^n}{(n-1)!}v_{n-1},\qquad v_0=1,
 \tag{GV.9}
\]
which is (GV.2). This argument specifies the quotient density; it
does not replace it by an unrelated quotient Riemannian metric.

For $J_{p,q}=\operatorname{diag}(1_p,-1_q)$, multiplication
$X\mapsto J_{p,q}X$ carries its unitary Lie algebra to
$\mathfrak u(n)$. Its complexification is multiplication by
$J_{p,q}$ on $M_n(\mathbb C)$, of complex determinant of modulus
one. It therefore preserves the exact Omega measure. An invariant
smooth Haar density is determined by its density at the identity;
the exponential map has derivative the identity there. The small-ball
limit in \texttt{OmegaDensity} is exactly that density relative to
$\lambda_n$, for the source's entry-sum balls as well as Euclidean
balls. The \texttt{CDHaar.prob} and \texttt{CDHaar.trick} fields now
force
\[
 \texttt{cd.cd}(p,q)\text{ to have Lie density }
                  A_n^{-1}\lambda_{p,q}.
 \tag{GV.10}
\]
In particular this holds for every supplied \texttt{CDHaar}, with
the same explicit constant $A_n^{-1}$ for every signature.

\paragraph{GV6: a rational gauge and its determinant under isometries.}
Choose $0\ne\delta\in E$ with $\bar\delta=-\delta$.
Use the following $F$-basis $\mathcal E$ of skew-Hermitian matrices:
$\delta E_{ii}$, and $E_{ij}-E_{ji}$,
$\delta(E_{ij}+E_{ji})$ for $i<j$.
On $\mathfrak u_J$ define a rational top form $\omega_J$ by declaring
the linear map $X\mapsto JX$ to identify its tangent density with
the coordinate form in $\mathcal E$. Translate this form on the group.
This tangent identification is not asserted to be a Lie homomorphism.

For any $B$ over a local quadratic algebra,
\[
 \det_F(Y\mapsto B^*YB)=N_{E/F}(\det B)^n
 \quad\text{on skew-Hermitian matrices}.
 \tag{GV.11}
\]
One proof is to split the quadratic algebra. The skew-Hermitian
matrices then identify with $M_n$, and the map is left and right
multiplication by the two component matrices of $B$, whose determinant
is the displayed norm power. Equality descends because these are
polynomial determinant identities. This proof covers a split place
and a dyadic place without dividing by two.

If $B^*JB=I$, the conjugation $g\mapsto B^{-1}gB$ transforms
$Y=JX$ into $B^*YB$. Its tangent determinant is $(\det J)^{-n}$.
Consequently, if it carries $K_{L_v}$ to $K_{\mathrm{std},v}$,
\[
 |\omega_J|_v(K_{L_v})
       =|\det J|_v^n\,|\omega_I|_v(K_{\mathrm{std},v}).
 \tag{GV.12}
\]
At a real place with $A^*J_{p,q}A=\tau J$, the corresponding map on
$Y$ is $Y\mapsto A^{-*}YA^{-1}$, again with absolute determinant
$|\det\tau J|^{-n}$. If $t_\tau=|\tau\delta|$, exactly $a_n$ of
the coordinates in $\mathcal E$ acquire a factor $t_\tau$.
By GV5,
\[
 |\omega_J|_\tau
   =|\det\tau J|^n\,t_\tau^{-a_n}v_n\,
                         \mu_{W,L,\tau}.
 \tag{GV.13}
\]
Here $\mu_{W,L,\tau}$ is the very real factor used in GV3.

\paragraph{GV7: the integral Lie lattice at dyadic as well as odd places.}
Let $R=\mathcal O_{F_v}$ and $S=\mathcal O_{E\otimes F_v}$.
An unramified quadratic algebra has an anti-invariant unit
$\kappa\in S$ with $\bar\kappa=-\kappa$. In the split case take
$(1,-1)$. In the nonsplit case lift an element whose residue
$z-\bar z$ is nonzero; such a residue exists also in characteristic
two because the trace map of the quadratic finite-field extension
is onto. Thus $S^- =R\kappa$.

The integral Lie lattice is the direct sum of $n$ copies of $S^-$
on the diagonal and $b_n$ copies of $S$ off the diagonal. The diagonal
coordinate determinant relative to $\delta$ has absolute value
$|\delta^2|_v^{-1/2}$. An integral basis of $S$ has unit
discriminant, whereas the discriminant of $(1,\delta)$ is
$4\delta^2$. The off-diagonal determinant relative to that basis
therefore has absolute value $|2|_v^{-1}|\delta^2|_v^{-1/2}$.
The determinant of the whole integral Lie lattice relative to
$\mathcal E$ is exactly
\[
 d_v=|2|_v^{-b_n}|\delta^2|_v^{-a_n/2}.
 \tag{GV.14}
\]
The factor $|2|_v^{-b_n}$ is essential; omitting it would give an
incorrect archimedean-to-finite normalization.

For clarity, the integral unitary group is smooth at every such
place by a direct lifting calculation. If a matrix is unitary modulo
$\varpi^r$, multiplication by $1+\varpi^r X$ changes its Hermitian
error modulo $\varpi^{r+1}$ by $X+X^*$. On off-diagonal entries this
map is onto by freely choosing one entry of each pair; on the
diagonal it is the surjective residue trace $S\to R$. Its kernel
has residue dimension $n^2$. Successive corrections converge in
the complete DVR. In particular reduction onto the finite unitary
group is surjective, and each subsequent congruence fiber has
$q_v^{n^2}$ elements. Equivalently, the same equations give local
analytic coordinates with unit Jacobian and $n^2$ free variables.
An invariant integral top density gives each first congruence
neighborhood mass $q_v^{-n^2}$. This supplies the measure assertion
in addition to a finite-set lifting assertion.

\paragraph{GV8: the exact finite normalized factors.}
For the standard self-dual lattice, GV7 gives
\[
 |\omega_I|_v(K_{\mathrm{std},v})
       =d_vq_v^{-n^2}|U_n(k_v)|
       =d_v\prod_{j=1}^n(1-\epsilon_v^j q_v^{-j}).
 \tag{GV.15}
\]
The notation $U_n(k_v)$ means $\mathrm{GL}_n(k_v)$ at a split
place. There its order follows by choosing independent columns.
At a nonsplit place, the number of norm-one vectors in $k_{v^2}^n$
is $q_v^{n-1}(q_v^n-(-1)^n)$: separate the last coordinate and use
the $q_v+1$ preimages of every nonzero norm, as in MF6. The unitary
group acts transitively on them, with stabilizer $U_{n-1}(k_v)$;
induction yields (GV.15). The residue-trace argument in GV7
establishes the same formula at $v\mid2$.

For a general self-dual $L_v$, DR1 gives an integral orthonormal
basis in the nonsplit case. At a split place choose a basis in one
factor and the conjugate inverse-transpose basis in the other.
Together these provide $B^*JB=I$ carrying the standard lattice to
$L_v$. Combining (GV.12) and (GV.15), relative to Haar measure
$dg_v$ with $K_{L_v}$ of mass one, gives
\[
 |\omega_J|_v
  =|\det J|_v^n d_v
                  \prod_{j=1}^n(1-\epsilon_v^j q_v^{-j})\,dg_v.
 \tag{GV.16}
\]
All determinants here are local coordinate determinants; no
globally free basis for $L$ has been introduced.

\paragraph{GV9: the global conversion constant, with no unspecified scalar.}
For $n>0$ define the adelic gauge normalization
\[
 d\tau_J=|\operatorname{disc}F|^{-n^2/2}L_F(1,\epsilon)^{-1}
        \prod_{\tau\mid\infty}|\omega_J|_\tau
        \prod_{v<\infty}L_v(1,\epsilon)|\omega_J|_v,
 \quad L_v(1,\epsilon)=(1-\epsilon_vq_v^{-1})^{-1}.
 \tag{GV.17}
\]
This is the convention of the classical theorem in GV10.
Its local products converge: the factor $L_v(1,\epsilon)$ cancels
the $j=1$ term of (GV.16), leaving $1+O_n(q_v^{-2})$; the
coefficient and determinant exceptions have finite support.

Multiply (GV.13) and (GV.16). First
$\prod_v|\det J|_v^n=1$ by the product formula in $F$.
Next $\delta^2\in F^\times$, $|\tau(\delta^2)|=t_\tau^2$, and
\[
 \left(\prod_{v<\infty}d_v\right)
                  \left(\prod_{\tau\mid\infty}t_\tau^{-a_n}v_n\right)
     =2^{b_nf}v_n^f=A_n^f.
 \tag{GV.18}
\]
Indeed the $\delta^2$ powers cancel by the product formula, and
$\prod_{v<\infty}|2|_v=2^{-f}$. The remaining finite Euler products
are absolutely convergent, so rearranging the finitely many $j$-products
is legitimate. We obtain the exact equality of measures
\[
 \boxed{\quad d\tau_J=C_n\,d\mu_{W,L},\qquad
 C_n=\frac{|\operatorname{disc}F|^{-n^2/2}A_n^f}
                   {\displaystyle\prod_{j=1}^nL_F(j,\epsilon^j)}.
       \quad}                                      \tag{GV.19}
\]
This is independent of $J$, the self-dual lattice, the rational
frame, the signature at every real place, and the chosen
$\texttt{CDHaar}$. Each independence follows from the calculated
factors. In particular (GV.19) does not assume signature independence
of quotient volumes or a mass formula.

\paragraph{GV10: the one retained Tamagawa theorem and finite covolume.}
We use the following established theorem as an explicit external input:
for every CM quadratic $E/F$ and every nonzero nondegenerate Hermitian
space, the measure (GV.17), using its usual local gauges and quadratic
$L$-value, has
\[
                    \tau_J\bigl(G_J(F)\backslash G_J(\mathbb A_F)\bigr)=2.
 \tag{GV.20}
\]
This includes finiteness. A precise published normalization and statement
are Harris--Li--Skinner, \emph{p-Adic L-Functions for Unitary Shimura
Varieties I}, Documenta Mathematica, Extra Volume Coates (2006),
\S0.2, p.~400, equation (0.2.1) and the following Tamagawa-number
statement; equation (0.2.2) gives its hyperspecial local convention.
The classical treatment is Weil, \emph{Adeles and Algebraic Groups},
Chapter IV, \S4.4.
\footnote{Primary publisher text:
\url{https://ems.press/content/book-chapter-files/27474};
DOI \url{https://doi.org/10.4171/DMS/4/12}.}
The theorem itself is not reconstructed here. It is stronger than
an assertion that an unspecified Haar measure has finite quotient
volume: its normalization is precisely (GV.17).

For an unramified extension with a self-dual lattice, GV9 therefore
gives, for every fundamental domain $D$, the positive finite value
\[
 \mu_{W,L}(D)=M_n,\qquad
 M_n=\frac{2|\operatorname{disc}F|^{n^2/2}}
                   {A_n^f}\prod_{j=1}^nL_F(j,\epsilon^j)
                   \quad(n>0).                     \tag{GV.21}
\]
For $n=0$ the group and rational subgroup are trivial and the
normalization in GV3 is the point mass one; put $M_0=1$.
The more general source statement \texttt{FiniteVolume\textsubscript{2}}
also follows from the same classical theorem, even when the CM
extension is ramified or no self-dual lattice is specified: any Haar
measure on this group is a positive finite scalar multiple of its
Tamagawa measure. GV4 supplies its fundamental domains and the common
mass. This general finiteness assertion does not extend (GV.19) to
ramified local normalizations.

\paragraph{GV11: the definite geometric quotient really is the genus mass.}
Assume $W$ is totally positive definite, and set
$K=G_J(F\otimes\mathbb R)K_L$. Its $\mu_{W,L}$-mass is one by
\texttt{CDHaar.prob} and (GV.6). The map
\[
 g\longmapsto L_g=W\cap g_f\widehat L
 \tag{GV.22}
\]
identifies $\Gamma\backslash G_J(\mathbb A_F)/K$ with rational
isometry classes in the genus of $L$. Here intersection and localization
are the actual operations of LP1. A tuple of local isometries is
integral outside finitely many places and so gives an adele;
conversely its intersection has exactly the prescribed localizations. At each of the finitely many exceptional
places, approximate a completed basis by rational vectors closely
enough to preserve its integral span; invertibility and the lattice
stabilizer are open. LP1 then applies to actual rational local modules.
The source genus uses rational local actions. MF3's Cayley weak
approximation places a rational unitary matrix in each needed open
local transporter, so the completed and rational-local formulations
of this genus agree. This does not assert strong approximation.

The positive trace-form proof of TL1 applies to this definite form as
well as to the coordinate identity form: self-duality bounds the
integral trace Gram determinant, positivity bounds multiplication
operators by $\max_\tau|\tau(a)|$, and the same short-vector argument
gives finitely many classes and finite automorphism groups. For a
representative $g_i$, its rational stabilizer in $g_iKg_i^{-1}$ is
exactly $\operatorname{Aut}(L_{g_i})$: preserve all completions and
intersect, or complete a global lattice automorphism. Denote its
order by $s_i$.

The normalization can now be checked without a volume theorem.
The nonnegative function
\[
                  \Phi(g)=\sum_i s_i^{-1}1_{g_iK}(g)
 \tag{GV.23}
\]
has $\sum_{\gamma\in\Gamma}\Phi(\gamma g)=1$ for every $g$.
Indeed exactly one double coset occurs, and in it the number of
translates meeting $g_iK$ is $s_i$. Integrate this equality on a
fundamental domain and use Tonelli as in GV4. Since each $g_iK$ has
measure one,
\[
 \boxed{\ \mu_{W,L}(D)=\sum_i\frac1{s_i}
                    =\operatorname{mass}(W,L).\ }   \tag{GV.24}
\]
The source's fallback definition of \texttt{mass} is harmless here:
the actual finite class representatives have just been supplied.
GV11, including finite covolume in the definite case, is independent
of the Tamagawa input in GV10.

\paragraph{GV12: exact lattice and rational-isometry transport.}
For two self-dual lattices in the same Hermitian space, DR1 and the
split-pair construction in GV8 give a finite adelic unitary element
$u$ carrying one completion to the other; take $u_v=1$ outside the
finite set where the lattices differ. Thus their compact open
stabilizers are conjugate. By GV1 a fixed finite Haar measure gives
them equal mass, so the normalized measure in GV3 is the same for
both lattices. This proves the literal \texttt{LatticeIndep}.

If a rational isometry $e:W_1\to W_2$ carries the specified lattices,
let $T$ be its matrix in the two specified rational frames.
Conjugation $\varphi(g)=TgT^{-1}$ is a topological group isomorphism
carrying the rational subgroup, the completed stabilizer, and the
real rectangle data to their counterparts. In particular
$\varphi(g)T=Tg$, exactly \texttt{IsConjBy}. Uniqueness in GV3 gives
$\varphi_*\mu_1=\mu_2$, and the partition definition of a fundamental
domain proves both directions of its transport clause. If $e$ does
not carry the specified lattices, first replace $L_2$ by $eL_1$ and
use lattice independence. Under the literal supplied \texttt{VolTie},
the equal domain masses then give $D.\mathrm{vol}(W_1)=D.\mathrm{vol}(W_2)$.
This proves all three clauses of \texttt{VolIndep} without Tamagawa
or a global isometry between different signatures.

\paragraph{GV13: actual signature independence and \texttt{VolumeEqMass\textsubscript{2}}.}
Now use the named input in GV10. Formula (GV.21) gives equal normalized
fundamental-domain masses for any two self-dual Hermitian spaces of
the same rank. This includes varying signatures at different real
places, and in particular all quantifiers in the source's
\texttt{VolSigIndep}. Apply GV11 to the actual positive standard
space $(E^n,I,\mathcal O_E^n)$. Its mass is the source's
$\mathrm{massI}(D,n)$, and hence
\[
          \mathrm{massI}(D,n)=M_n,\qquad
          \mu_{W,L}(D_W)=\mathrm{massI}(D,n).
 \tag{GV.25}
\]
Thus every actual positive self-dual genus has the same absolute mass;
no prior identification of rational Hermitian spaces was used.

For a supplied record $D$ with \texttt{VolTie}, its own normalized
domain mass is $D.\mathrm{vol}(W)$. Positivity permits passing from
$\mathrm{ofReal}$ to real numbers, so
\[
              D.\mathrm{vol}(W)=\mathrm{massI}(D,\operatorname{rank}W)
 \tag{GV.26}
\]
for every constant-signature $W$. The signature sum gives
$\operatorname{rank}W=P+Q$. This proves both literal clauses of
\texttt{VolumeEqMass\textsubscript{2}}, including the positive clause
in rank zero and the indefinite clause with $2\le P$.
The proof actually gives the indefinite equality without that lower
bound. It does not assign this volume to an arbitrary preexisting
positive function called $D.\mathrm{vol}$; its \texttt{VolTie}
hypothesis is essential.

\paragraph{GV14: a norm-sign vector and a self-dual orthogonal complement.}
This paragraph uses the named Hasse--Minkowski theorem retained in MF4.
Let $W$ have rank at least three, constant signature $(P,Q)$, and
choose $s=1$ with $P\ge1$ or $s=-1$ with $Q\ge1$.
At each finite place a local orthonormal basis exists by GV8 and the
unit norm map is onto by DR1, so $h(x,x)=s$ has a local solution.
At real places it has a solution by the stipulated signature.
On the underlying $F$-space let $q(x)=h(x,x)$.
The quadratic form $q\perp\langle-s\rangle$ is locally isotropic
everywhere. Hasse--Minkowski gives a global isotropic pair $(x,t)$.
If $t\ne0$, replace $x$ by $x/t$. If $t=0$, $q$ has a nonzero
isotropic vector; nondegeneracy produces a hyperbolic plane, which
represents every scalar, in particular $s$. Thus there is a rational
$x$ with $h(x,x)=s$.

Let $W'=x^\perp$. The formula
$y\mapsto y-s^{-1}x h(x,y)$ is the rational orthogonal projection,
so $W=Ex\perp W'$ and $W'$ has the required decremented signature.
We must supply its lattice, rather than assume that $L\cap W'$
is self-dual at every place. At a finite place choose a unit of norm
$s$. MF3's equal-norm reflection transports its first coordinate
vector to $x$, leaving a standard self-dual lattice on the
orthogonal complement. At a split place the same construction is
elementary linear algebra: choose a first vector and a complementary
kernel of its paired covector, then use the inverse-transpose dual
basis in the other factor. At all but finitely many places $x\in L_v$;
there the projection is integral, so
$L_v=\mathcal O_{E_v}x\perp(L_v\cap W'_v)$ and its complement is
already self-dual. Approximate their finitely many exceptional completed bases by rational
bases inside the same open lattice transporter, then patch the resulting
rational local complement lattices by LP1. The result is a full global self-dual lattice in
$W'$. Choose any rational frame for $W'$ and restrict the form.
Its inclusion into $W$ has range $\ker h(x,-)$ and is isometric.
These are precisely all the \texttt{UnitSplit} fields. A global
integral orthonormal basis and a global Landherr normal form were
not required.

\paragraph{GV15: the geometric successor ratio and \texttt{Conclusion2}.}
For $n\ge3$, divide (GV.21) in ranks $n$ and $n-1$:
\[
 \frac{M_n}{M_{n-1}}
  =|\operatorname{disc}F|^{n-1/2}
       \left(\frac{(n-1)!}{(2\pi)^n}\right)^f
                     L_F(n,\epsilon^n).
 \tag{GV.27}
\]
MF1 proves $f=\mathrm{fDeg}(E)$,
$|\operatorname{disc}F|=(2\,\mathrm{vInv}(E))^f$, and
$L_F(n,\epsilon^n)=\mathrm{lValEta}(E,n)$, with the two split-prime
half powers paired exactly. Substitution turns (GV.27) into
\[
 \frac{\big((n-1)!(\mathrm{vInv}(E)/\pi)^n\big)^f}
            {\sqrt{(2\,\mathrm{vInv}(E))^f}}
                       \mathrm{lValEta}(E,n),         \tag{GV.28}
\]
the literal \texttt{massRatioRightSide}. All divisors are positive.
For the actual $x,W',\iota$ of GV14, GV13 and \texttt{VolTie}
identify the left side with $D.\mathrm{vol}(W)/D.\mathrm{vol}(W')$.
This proves every clause of \texttt{Conclusion2}, and, before using
\texttt{VolTie}, its actual-normalized-domain version
\texttt{VolRatioTrue}. It is consistent with MF19, but (GV.27) was
derived from the geometric quotient, not by relabeling MF19's
mass law as a volume law. Its classical premises are exactly those
stated in GV10 and GV14.

\paragraph{GV16: the source's chosen \texttt{trueVol} has its intended value.}
For every $W$ satisfying \texttt{SigOK}, GV3 gives at least one
normalized measure and proves the property for every such measure.
GV4 gives at least one fundamental domain and proves equality for
every such domain. GV10 gives their common positive finite real
value $M_n$. Consequently
\[
 \mathrm{VolIs}(\mathrm{cm},\tau,\mathrm{cd},W,M_n),\qquad
 \mathrm{trueVol}(\mathrm{cm},\tau,\mathrm{cd},W)=M_n.
 \tag{GV.29}
\]
For the second equality, the defining existential branch is now
inhabited; its chosen value has the same domain mass and is positive,
so uniqueness forces equality. The fallback value $1$ is never used
on these valid spaces. This proves the actual
\texttt{volIs\_trueVol} supplier and hence
\[
 [\,D.\mathrm{vol}(W)=\mathrm{trueVol}(D.\mathrm{sp.cm},
              \tau_D,D.\mathrm{cd},W)\text{ for every }W\,]
                 \ \Longrightarrow\ \mathrm{VolTieG}(D).
 \tag{GV.30}
\]
The host CM-type is the one in the original record; its bijection
with infinite places is part of the host interface, not a change
of embeddings. Constant-signature specialization also gives
\texttt{VolTie}. For malformed signature fields outside
\texttt{SigOK}, no assertion that a chosen fallback has geometric
meaning is made.

\paragraph{GV17: the exact compact-dual scale in every sign ordering.}
Fix a Boolean sign function $e:\{1,\ldots,n\}\to\{+,-\}$ and
a sorting bijection $\rho:\{1,\ldots,P\}\sqcup\{1,\ldots,Q\}
\to\{1,\ldots,n\}$. Reindexing matrices is a topological group
isomorphism $U(P,Q)\to U(J_e)$. Its pushforward of
$\texttt{cd.cd}(P,Q)$ and the source's fixed
$\mathrm{REAL.muG}(e)$ are Haar measures, so there is a unique
positive finite scalar $c_e$ relating them. If a different sorting
is chosen, the two maps differ by conjugation with a permutation
unitary matrix; GV1 makes the scalar unchanged. Thus the same $c_e$
works for every sorting and every measurable set, exactly as demanded
by \texttt{CdScaleIs}. Positivity and uniqueness can be checked on
any compact set with nonempty interior.

Choosing these $c_e$ for all $n,e$ inhabits the existential branch
in \texttt{trueScale}; uniqueness proves
\[
               \mathrm{CdScaleIs}(\mathrm{cd},
                                  \mathrm{trueScale}(\mathrm{cd})).
 \tag{GV.31}
\]
Here no numerical normalization of \texttt{REAL.muG} is silently
assumed: $c_e$ is its exact positive Haar ratio, and the defining
identity on all measurable sets is proved. The explicit numerical
normalization of \texttt{cd.cd} itself was already computed in GV5.
This proves \texttt{cdScaleTie} whenever the same record's scale
function equals \texttt{trueScale}. It is independent of the
Tamagawa theorem and of any discrete-series existence statement.

\paragraph{GV18: the actual quotient probability measure and level changes.}
Push the restriction of $\mu_{W,L}$ to a fundamental domain onto
$\Gamma\backslash G_J(\mathbb A_F)$. For a nonnegative measurable
function on this quotient, its pullback is $\Gamma$-invariant;
the argument of (GV.7) shows that the integral is independent of
the selected domain. Thus the quotient measure, not only its total
mass, is canonical. Its mass is $M_n$, and its probability
normalization is precisely
\[
                       \mu_Q^1=M_n^{-1}\mu_Q.
 \tag{GV.32}
\]
If the same data have \texttt{VolTieG}, replace $M_n$ by
$D.\mathrm{vol}(W)$ in (GV.32). This proves the source's
\texttt{VolPairEx}, \texttt{FinVol},
\texttt{normalized\_quotient\_measure\_eq}, and
\texttt{probability\_measure\_eq\_pair}; none of their fallback
branches is used on a valid space.

There is also an exact finite-level check. If $K'\subseteq K_L$
is an open subgroup of finite index $r$, then
$\mu_f(K')=1/r$. Renormalizing the finite Haar measure to give
$K'$ mass one multiplies the adelic quotient mass by $r$ and leaves
the normalized probability measure unchanged. More generally the
ratio for two compact open levels is the ratio of their indices
over their common intersection. This follows from their finite
coset decompositions, and fixes the direction of every level factor.
It does not identify an unspecified model level with $K_L$.

\paragraph{GV19: the original base-volume stage, on the same base and columns.}
IG12 constructs, for the actual $B,C$ and supplied
$\mathrm{cd}$, a stage $G$ with $G.\mathrm{cd}=\mathrm{cd}$ and
\[
 G.\mathrm{vol}(W)=\mathrm{baseVol}(B,G.\mathrm{cd},W)
  =\mathrm{trueVol}(B.\mathrm{sp.cm},
          (\tau\mapsto B.\mathrm{host.emb}(\tau)),G.\mathrm{cd},W).
 \tag{GV.33}
\]
GV16 now upgrades this exact original stage from positivity to
actual \texttt{VolTieG}. For every choice of the other stages
$W_0,V,S,N$ over this same $B,C$, the assembled record has the
same host, CM-type, Hermitian spaces, Haar family and volume
function. Therefore (GV.30) holds for every one of those assembled
records. This is the universal volume clause used in
\texttt{GIs2}; it does not require changing the row or local stages.

The same argument applies to any $G.\mathrm{withTrue}$: by its
literal definition only \texttt{vol}, its positivity proof, and
\texttt{gx} are reset, while $B,C,\mathrm{cd}$ and the jet data
remain fixed. A further reset of \texttt{gat} and its size proof
also leaves (GV.33) untouched. Thus the volume clause in
\texttt{exists\_trueVolume\_stage\_of\_any\_column\_choice} is
now supplied. Its other \texttt{ColumnDataSplitPrimes} and
$\Omega$-compatibility clauses still require their exact standard
column, local-family and coefficient data. A dependent layout or
Core on a previous record is not automatically transported across
\texttt{withTrue}. Empty jets in IG12 remain empty, and no nonempty
arithmetic twist is created by this upgrade.

\paragraph{GV20: precisely which real-data fields have been supplied.}
The literal \texttt{RealTied}(D) has four fields:
$\mathrm{cdScale}$, $\mathrm{dsig}$, $\mathrm{dsig\_all}$,
and $\mathrm{vol}$. For the same record $D$, assume
\[
 D.s.\mathrm{cdScale}=\mathrm{trueScale}(D.\mathrm{cd}),\qquad
 D.\mathrm{dsig}(\tau)=\mathrm{trueDsig}(D.R,D.\mathrm{hc}(\tau)).
 \tag{GV.34}
\]
GV17 supplies the first field, and GV16 or GV19 supplies the last
field when $D$ has the actual base-volume normalization.

For the middle fields we retain exactly the source's
\texttt{DiscreteSeriesExist} and \texttt{FormalDegreeFormula\textsubscript{2}}.
The former supplies a discrete-series representation for each
regular Harish--Chandra parameter; the latter identifies its
formal degree, for the Haar measure in (GV.31), with the displayed
Harish--Chandra product. With the supplied parameter laws,
\texttt{trueDsig} is precisely that product. One chosen representation
therefore gives \texttt{IsDsig}; applying the same formula to every
representation and every nonzero vector gives \texttt{dsigTie}.
This proves the exact conditional \texttt{RealOf} join, with
(GV.34) imposed on this $D$. It does not derive representation
existence or a formal-degree formula from the Tamagawa theorem,
and it does not treat the larger \texttt{retie} operation as the
identity: \texttt{retie} also changes \texttt{dsig}, scale and
\texttt{gx}, with their dependent consequences.

\paragraph{GV21: the honest same-object \texttt{V4.Core} consequence.}
Fix the actual $K,X,f,T$, put $D=X.\mathrm{toD7}$, and suppose
there is a layout $Y:\mathrm{Lay}(D)$ with the hypotheses of
\texttt{coreOfLay}. Impose (GV.33)--(GV.34) on this very $X$.
Then GV20 supplies its \texttt{RealTied} argument, and GV13 supplies
the geometric \texttt{VolumeEqMass\textsubscript{2}} consequence
used by the size and formal-degree joins. The constructor keeps
$O=\mathrm{layObj}(Y)$ and all its levels, lattices, comparisons,
theta objects and arithmetic choice. Its fields are verified as
follows: \texttt{RecOK} supplies the column bounds and the
\texttt{Del}/\texttt{etaF} identities; \texttt{LedOK}, together with
$\mathrm{MuIsW.off}$, supplies \texttt{fval} and the actual Gy
bounds; $Y$ supplies the remaining column, theta, gauge and
$\mathrm{NChTie}$ fields; the exact untwisted vanishing clause
completes \texttt{RowLay}; and GV20 supplies \texttt{real}.

The other arguments remain explicit: \texttt{Over}, the same-map
$\xi$ and \texttt{MuIsW}, the forced-prime inclusion,
\texttt{GatNum}, untwisted row vanishing, \texttt{MtStableAt},
the primitive type-character powers, and
\[
   Y.\mathrm{nch}.Q(w)=T.Q(w),\qquad
   Y.\mathrm{nch}.k'(w)=T.k(w)\qquad(w\in W_t).
 \tag{GV.35}
\]
These last identities are conductor-support and exponent identities,
not merely a choice of a prime or a side. This is a strengthened
conditional Core supplier with the geometric volume and scale
arguments proved; it is not a proof that an arbitrary original
stage admits all the other arguments. In particular the forced-prime
field contains the distinguished prime $g_0$, so an empty twist-prime
set cannot satisfy it. The bare empty models of earlier stage
existence results still do not produce a \texttt{V4.Core} or
\texttt{ExTrue}. The conditional Tamagawa and formal-degree inputs
are visible at the final endpoint, and no RH or GRH conclusion is
drawn from these geometric identities alone.
